% Generated draft: requires mathematical and source review.
\documentclass[11pt]{article}
\usepackage[margin=25mm]{geometry}
\usepackage{fontspec}
\setmainfont{Latin Modern Roman}
\usepackage{amsmath,amssymb,mathtools,unicode-math}
\setmathfont{Latin Modern Math}
\usepackage{enumitem,needspace,xurl,hyperref}
\setlength{\parindent}{0pt}
\setlength{\parskip}{5pt}
\setlength{\emergencystretch}{3em}
\newcommand{\N}{\mathbb{N}}
\newcommand{\Z}{\mathbb{Z}}
\newcommand{\Q}{\mathbb{Q}}
\newcommand{\R}{\mathbb{R}}
\newcommand{\C}{\mathbb{C}}
\newcommand{\F}{\mathbb{F}}
\newcommand{\A}{\mathbb{A}}
\newcommand{\PP}{\mathbb{P}}
\newcommand{\E}{\mathbb{E}}
\newcommand{\eps}{\varepsilon}
\newcommand{\dd}{\,\mathrm{d}}
\begin{document}
\section*{UnsolvedMath ID: 30005944}
\textbf{Source problem code:} catalogue entry
\subsection{Definitions and mathematical statement}
For a smooth proper curve \(C\) over a field \(k\), a linear series of degree \(d\) and dimension \(r\) is a pair \((L,V)\), where \(L\) is a line bundle on \(C\) with \(\deg L=d\), and \(V\subseteq H^0(C,L)\) is a \(k\)-linear subspace of dimension \(r+1\). Here \(d\) and \(r\) describe the numerical data of a linear series; the supplied question does not fix them, specify the ground field or its characteristic, or assume completeness or absence of base points.

\textbf{Problem.} Given a degeneration of a smooth curve \(C\) to a nodal curve \(C_0\), and a linear series \((L,V)\) on the smooth curve, develop a description of the limiting data and their moduli that addresses the following questions:
\begin{enumerate}
\item How does the line bundle \(L\) degenerate, and how does the subspace \(V\) degenerate together with it?
\item What is the appropriate moduli space of linear series on \(C_0\), and what can be said about that space?
\item For a linear series on \(C_0\), what conditions characterize its occurrence as the limit of linear series on smooth curves?
\end{enumerate}
In the third question, occurrence as a limit requires a family with nodal special fibre \(C_0\), smooth general fibre, and compatible linear-series data specializing to the proposed data on \(C_0\). The meaning of compatibility and specialization depends on the theory of linear series adopted on nodal curves. The source does not specify that theory, whether the smoothing family is fixed or may vary, or whether base change is permitted. Thus these are questions seeking a moduli and degeneration theory, rather than a fully specified representability assertion or a single conjectural equivalence. No restriction to curves of compact type is stated.

\textbf{Supplied logarithmic context.} The accompanying excerpt defines a log curve over \(S\) as a proper, integral, logarithmically smooth morphism \(\pi\colon C\to S\) with connected, one-dimensional geometric fibres. It describes an associated morphism of Olsson fans \(\mathcal A_C\to\mathcal A_S\) and the relative object
\[
\Gamma=S\times_{\mathcal A_S}\mathcal A_C.
\]
The excerpt presents this object as combinatorial through a stated correspondence between Olsson fans and generalized cone complexes. It attributes a classification of log smooth curves to F.\ Kato, with reference label [6], and uses reference labels [1] and [2] for the fan constructions and correspondence. Those reference entries are not supplied. The excerpt then refers to a tropicalization \(\tau\) without supplying its definition and ends mid-sentence. Consequently it does not determine the logarithmic linear-series objects or moduli functor needed to make the three questions more specific.
\subsection{Short English statement}
Describe how a line bundle and its chosen space of sections behave when a smooth curve acquires nodes, construct suitable moduli of the resulting linear series, and characterize which such objects can be smoothed. The supplied record leaves the ground field, numerical invariants, degeneration setup, and precise notion of a linear series on a nodal curve unspecified. It reports partial progress, but this draft does not independently verify that assessment.
\subsection{Sources}
\begin{itemize}
\item Supplied local UnsolvedMath registry record, permanent ID \(30005944\), source problem code \texttt{catalogue entry}, titled \emph{Degeneration and Moduli of Logarithmic Linear Series}. Project registry locations: \texttt{data/unsolvedmath/problems.json} and supplementary \texttt{catalog/problems.json}. The supplied record is the local evidence used here; these files were not independently inspected.
\item \emph{Real and Logarithmic Enumerative Geometry} (2024), Oberwolfach Reports, DOI 10.4171/owr/2024/23, pp.\ 1285--1338. Supplied report link and literature-source label: \emph{Oberwolfach Report 23/2024}, \url{https://doi.org/10.4171/owr/2024/23}. The excerpt identifies Question 1 in \emph{Logarithmic linear series}, attributed there to Franchesca Carocci, joint work with Luca Battistella and Jonathan Wise. The registry supplies no separate proposer.
\item Supplied literature source: \emph{Vector bundles on Olsson fans}, \url{https://arxiv.org/abs/2602.05685}. No authors or more precise supporting passages for this reference are supplied in the record.
\item Literature assessment dated exactly \texttt{2026-08-22}: the supplied curation labels the problem \texttt{partially\_solved}. It states that logarithmic linear-series constructions address degeneration in important cases and extend limit-linear-series methods beyond compact type, while the 2026 foundational work still describes the full moduli problem as developing. This is supplied dated curation, not an independently verified literature assessment; the record itself requests independent verification of time-sensitive status and cited claims.
\item The registry labels the statement \texttt{corrected\_verified} and reports that it was checked against the official Oberwolfach report. This is a supplied source-verification claim, not verification performed for this draft. The report excerpt's references [1], [2], and [6] have no supplied bibliographic entries.
\end{itemize}
\end{document}
